%%=============================================================================
%%    Conceptuele Oplossing
%%=============================================================================

%    Startend vanuit de inleiding in het vorige hoofdstuk plaats je de mogelijke methodes en alternatieven naast
%    elkaar en ga je conceptueel na wat de mogelijke oplossingsmethodes kunnen zijn; daartoe kan je boeken,
%    websites en artikels raadplegen.

%    Dit hoofdstuk kan volgende zaken omvatten:
%    ● studie van normen en protocollen
%    ● studie van de werking en/of bediening van een machine of meettoestel
%    ● studie van een bestaand stappenplan (vb. bij onderhoud)
%    ● studie van een bestaande proefopzet of meetopstelling
%    ● studie en keuze van de juiste elektronische componenten
%    ● studie van de huidige situatie van een computernetwerk
%    ● functioneel ontwerp (eventueel met screenshots) van de gebruikersinterface van een programma of
%       website
%    ● studie en keuze van het besturingssysteem
%    ● studie en keuze van softwarepakketten
%    ● studie van het gekozen softwarepakket – gekozen hardware.

%    Je mag er van uit gaan dat alles wat je gezien hebt tijdens je opleiding gekend is bij de lezer. Nieuwe
%    technologieën en methodieken die je nodig hebt om je bachelorproefopdracht uit te voeren, worden
%    uiteraard wel in dit conceptueel hoofdstuk opgenomen en eventueel uitgewerkt in het relevante volgende
%    hoofdstuk.

%    Merk op: je hebt ongetwijfeld heel wat voor- en nadelen tegenover elkaar moeten afwegen om de juiste
%    keuzes te maken. Dit is héél interessant om te vermelden!
%    Op het einde maak je een keuze welke methode je gaat gebruiken. Deze oplossing wordt in de volgende
%    hoofdstukken uitgewerkt. Het is een goed idee om dit hoofdstuk een overzicht op te nemen welk deelaspect
%    in welk van de volgende hoofdstukken zal uitgewerkt worden. Je mag daarbij de hoofstuktitels op niveau 1
%    van je inhoudsopgave vermelden met een korte inhoud.
%    Zorg dat je voorstellen SMART zijn: Specifiek, Meetbaar, Aanvaardbaar, Realistisch en Tijdsgebonden.
%    Meer informatie over het hoofdstuk “Conceptuele oplossing” is te vinden in “Rapporteren voor technici” pg.
%    86 tot 87

\chapter{\IfLanguageName{dutch}{Conceptuele Oplossing}{}}%
\label{ch:conceptuele-oplossing}

Startend vanuit de inleiding in het vorige hoofdstuk plaats je de mogelijke methodes en alternatieven naast
elkaar en ga je conceptueel na wat de mogelijke oplossingsmethodes kunnen zijn; daartoe kan je boeken,
websites en artikels raadplegen.

Dit hoofdstuk kan volgende zaken omvatten:
\begin{itemize}[]
    \item studie van normen en protocollen
    \item studie van de werking en/of bediening van een machine of meettoestel
    \item studie van een bestaand stappenplan (vb. bij onderhoud)
    \item studie van een bestaande proefopzet of meetopstelling
    \item studie en keuze van de juiste elektronische componenten
    \item studie van de huidige situatie van een computernetwerk
    \item functioneel ontwerp (eventueel met screenshots) van de gebruikersinterface van een programma of website
    \item studie en keuze van het besturingssysteem
    \item studie en keuze van softwarepakketten
    \item studie van het gekozen softwarepakket – gekozen hardware.
\end{itemize}

Je mag er van uit gaan dat alles wat je gezien hebt tijdens je opleiding gekend is bij de lezer. Nieuwe
technologieën en methodieken die je nodig hebt om je bachelorproefopdracht uit te voeren, worden
uiteraard wel in dit conceptueel hoofdstuk opgenomen en eventueel uitgewerkt in het relevante volgende
hoofdstuk.

Merk op: je hebt ongetwijfeld heel wat voor- en nadelen tegenover elkaar moeten afwegen om de juiste
keuzes te maken. Dit is héél interessant om te vermelden!
Op het einde maak je een keuze welke methode je gaat gebruiken. Deze oplossing wordt in de volgende
hoofdstukken uitgewerkt. Het is een goed idee om dit hoofdstuk een overzicht op te nemen welk deelaspect
in welk van de volgende hoofdstukken zal uitgewerkt worden. Je mag daarbij de hoofstuktitels op niveau 1
van je inhoudsopgave vermelden met een korte inhoud.
Zorg dat je voorstellen \acrshort{SMART} zijn: \acrlong{SMART}.
Meer informatie over het hoofdstuk “Conceptuele oplossing” is te vinden in “Rapporteren voor technici” pg.
86 tot 87
